\section{Proposed Grid-Aware Demand Clustering for Location-Capacity Optimisation Methodology}

The proposed Grid-aware Demand Clustering for Location-Capacity Optimisation (GDC-LCO) framework consists of demand feature construction, grid-risk proxy construction, service-zone clustering, candidate station generation, and location-capacity optimisation.

\begin{figure*}[!t]
\centering
\begin{tikzpicture}[
node distance=6mm,
box/.style={draw, rounded corners, align=center, minimum width=21mm, minimum height=11mm, font=\scriptsize},
arrow/.style={-Latex, thick}
]
\node[box] (data) {Open urban\\data};
\node[box, right=of data] (features) {Demand, access,\\grid-risk features};
\node[box, right=of features] (cluster) {Grid-aware\\demand clustering};
\node[box, right=of cluster] (candidate) {Candidate site\\generation};
\node[box, right=of candidate] (optimise) {Location-capacity\\optimisation};
\node[box, right=of optimise] (eval) {Evaluation and\\visualisation};
\draw[arrow] (data) -- (features);
\draw[arrow] (features) -- (cluster);
\draw[arrow] (cluster) -- (candidate);
\draw[arrow] (candidate) -- (optimise);
\draw[arrow] (optimise) -- (eval);
\end{tikzpicture}

\caption{Workflow of the proposed GDC-LCO framework. Urban open data are transformed into demand, accessibility, and grid-risk features. These features are used to form service zones before candidate site selection, station sizing, and evaluation.}
\label{fig:method_flowchart}
\end{figure*}

\subsection{Demand Feature Construction}
Charging demand is represented as a weighted combination of residential, commercial, and road-related components:
\begin{equation}
D_i = \alpha D_i^{\text{res}} + \beta D_i^{\text{com}} + \gamma D_i^{\text{road}}.
\end{equation}

\subsection{Grid-Risk Proxy Construction}
In the absence of full distribution network topology, grid impact is represented by reproducible proxy variables such as peak-load pressure, assumed hosting capacity, existing charging density, or distance to an assumed connection point. A simple proxy can be expressed as:
\begin{equation}
G_i = \frac{\text{peak load proxy}_i}{\text{hosting capacity proxy}_i}.
\end{equation}

\subsection{Clustering Methods}
Three clustering methods are compared. Geographical K-means uses only spatial coordinates. Demand-weighted K-means uses demand weights when updating cluster centroids. The proposed grid-aware demand clustering uses:
\begin{equation}
z_i = [\lambda_s lon_i, \lambda_s lat_i, \lambda_d D_i, \lambda_a A_i, \lambda_g G_i].
\end{equation}

\subsection{Location-Capacity Optimisation Within Clusters}
After clustering, each service zone is optimised independently. This decomposition reduces the effective problem size and produces interpretable local planning regions. The current research code implements a lightweight greedy solver to maintain reproducibility without proprietary optimisation software, while the documented formulation supports later replacement with a mixed-integer or multi-objective optimiser.

\subsection{Algorithm Summary}
\begin{enumerate}
  \item Construct demand, accessibility, and grid-risk features.
  \item Partition demand points using grid-aware clustering.
  \item Generate candidate stations within each service zone.
  \item Select stations and charger capacities inside each zone.
  \item Aggregate selected stations and evaluate distance, cost, coverage, grid impact, and utilisation.
\end{enumerate}
